\scanpage{26}
g) $X$ muni d'une topologie $T$ déduite d'un préordre $\omega$ [donc les ouverts
sont les réunions d'ens de la forme $O_x$, $x\in X$]. $X$ est irréductible
sss deux ouverts non vides ont une intersection non vide, i.e. pour
$x,y\in X$, $O_x\cap\unsure{U_y|O_y}\ne\varnothing$ i.e. $\forall x,y\in X$, $\exists z\in X$ avec $z\leq x,\ z\leq y$,
i.e. $X$ est \uline{cofiltrant}. En conclure que les fermés irréductibles
de $X$ sont les parties $Y$ de $X$ telles que

\[
\begin{aligned}
&\text{a)}\quad y\in Y,\ \inserted{y'\in X,}\ y'\geq y\Rightarrow y'\in Y\qquad(Y\text{ fermé})\\
&\text{b)}\quad\forall x,y\in Y,\ \exists z\in Y,\ \text{tel que }\ z\leq x,\ z\leq y.
\end{aligned}
\]

h) $X$ comme dans g), $\underline X$ la catégorie associée à la relation de \inserted{pré}ordre
$\operatorname{card}\operatorname{Hom}_{\underline X}(x,y)\leq1$, $\deleted{\unsure{\mathrm{Hom}}}\operatorname{Hom}_{X}(x,y)\ne\varnothing\Leftrightarrow x\leq y\Longleftrightarrow y\in\overline{\{x\}}$ ---
Considérons le foncteur $\operatorname{Top}(X)\mathbin{\unsure{-|\to}}\widehat{\underline X}$, $F\longmapsto(x\longmapsto F(O_x))$.

Prouver que ce foncteur est une équivalence

\[
\operatorname{Top}(X)\mathrel{\overset{\sim}{\simeq}}\widehat{\underline X}
\]

Désormais on identifie [$\deleted{\unsure{X|\text{si}}}$ $X$ présobre i.e. \uline{ordonné}] $X$ à un
ensemble de sous-objets de l'objet final de $\widehat{\underline X}=\operatorname{Top}(X)$
[\deleted{\unsure{c.à.d|cribles}} d'ouverts du topos), et de $=$ pour les ouverts de $X$,
identifiés à des s-objets de $e_{\widehat{\underline X}}$ donc à des ouverts du topos.

Interpréter

\[
\underline{\operatorname{Points}}(\widehat{\underline X})\simeq\operatorname{Pro}(\underline X)
\]

en identifiant les pro-objets de $\underline X$ \inserted{(mod iso)} aux sous-ensembles
ordonnés $Y$ de $X$ satisfaisant les conditions a) b) de g),
\deleted{\unsure{avec|et}} la structure de catégorie provenant de la relation
d'ordre entre parties de $X$ \uline{opposée} de la relation d'inclusion.
Soit \deleted{\unsure{Privé de|$\operatorname{Pro}(X)$}}\inserted{$\widetilde X$} nouvel ens. ordonné ainsi construit. L'application
$x\longmapsto\overline{\{x\}}=\{y\in X\mid y\geq x\}$ est une application $X\mathbin{\unsure{-|\to}}\widetilde X$, qui
s'identifie à l'application can. de $X$ dans l'espace sobre
associé, \uline{ou} à $\underline X\longrightarrow\operatorname{Pro}(\underline X)$ (foncteur canonique).

i) \uline{Conditions équivalentes} $\alpha)$ $X\longrightarrow\widetilde X$ surjective
$\beta)$ Toute partie fermée cofiltrante de $X$ a un plus petit élément.
% archive sheet number: 25
